\section{One strategy against every continuation}\label{sec:strategy}

We now turn the final card into a single strategy. Making its choices
definite also clarifies that the height counts actual Maker moves.

\begin{proof}[Proof of Theorem~\ref{thm:main}]
Fix an enumeration of $\Z^2$: order cells by increasing
$\max(|x|,|y|)$, then lexicographically within each finite level. Retain the printed order of the children in every
combination expression. These choices are fixed before any play.
Initially the active claim is card $727$ in its identity placement.
Its hypotheses hold on the empty board by
\eqref{eq:final-card}.

At an active combination with current placement $F$, claim its
placed pivot $F(p)$ if free. If Maker already owns $F(p)$, claim
the first free cell in the fixed enumeration instead. The invariant
of Definition~\ref{def:claim} guarantees that Breaker does not own
the pivot, and finiteness of ownership guarantees a free replacement.
Stop if a target is complete. Otherwise, after Breaker's next legal
reply, choose the first child in the printed list whose placed
envelope avoids that reply. Lemma~\ref{lem:combination} proves that
such a child exists and that all its requirements and its full
envelope satisfy the invariant, including the older Breaker cells.
For a referenced card whose reference placement is $G$, the new
placement is $F\circ G$; for an inline expression, retain $F$ and
use that expression.

At a placed base card, either the target is already present or its
one missing cell is free. In the latter case claim it and win.
Thus a nonterminal step spends one actual Maker move and passes,
after one actual Breaker reply, to a child of strictly smaller
height. The bases have height $1$. Starting at height $21$, this
procedure wins within $21$ Maker moves.

The active expression and placement are determined by the complete
history: replay the prescribed pivot or replacement choices and
the first-surviving-child choices from card $727$.
On a history inconsistent with an earlier prescribed Maker choice,
use the first free cell of the enumeration. Use the same fallback
on an already-won history if one chooses to continue it formally.
This defines a globally legal policy on all finite legal histories,
including histories not reached from the empty board under the
policy. No choice of the policy depends on future Breaker moves.
The preceding induction applies to every legal continuation from
the empty board and proves the theorem.
\end{proof}

Only replies before victory are used in this argument. If the first
win occurs on Maker move $m\le21$, there have been precisely $m-1$
Breaker moves. In particular, no reply after Maker's winning move
is assumed or required.

One can express the same quantifier order using a fixed endpoint.
Extend the policy beyond a completed target by the first-free-cell
rule. For every sequence of proposed Breaker cells whose first
$20$ entries are legal along this extended policy, the Maker set
immediately after move $21$ contains an allowed copy of $S$.
An earlier completed copy persists because ownership is permanent.
This formulation is $\exists\sigma\,\forall\beta$, with $\sigma$
chosen in the proof and $\beta$ the sequence of Breaker replies;
it makes no assumption about a twenty-first Breaker reply.
Stopping at the first target recovers the ordinary game.

\begin{corollary}[A finite board]\label{cor:finite-board}
Maker wins the weak Maker--Breaker game within $21$ actual Maker
claims on the initially empty $251$-cell board $T_{727}$, and also
on the initially empty $17\times17$ board $\{0,\ldots,16\}^2$.
In each game, Maker moves first, the players alternate claiming one
previously unclaimed cell, and the targets are the allowed copies
of $S$ contained in the board.
\end{corollary}

\begin{proof}
Evaluation of the literal certificate gives
\[
 T_{727}\subseteq\{0,\ldots,16\}^2,
\]
in addition to the endpoint data in \eqref{eq:final-card}.
Every placed child envelope is contained in its parent envelope,
and every pivot and base target lies in its own envelope. Thus all
prescribed pivots and base completions lie in $T_{727}$.
Use that strategy, replacing an already-owned pivot by the first
free cell in a fixed ordering of $T_{727}$. Before Maker claim
$m\le21$, at most $2(m-1)\le40$ cells have been occupied, so such
a replacement is always available on either board.

After a pivot or replacement claim, Maker owns every child
requirement and the intersection of the child envelopes. Every
legal Breaker reply therefore leaves a surviving child envelope;
nesting also keeps the older Breaker cells outside that envelope.
This includes replies outside $T_{727}$ on the square board.
The same height induction now gives victory within $21$ actual
Maker claims, counting each replacement. Play stops as soon as a
target is complete; no Breaker reply after victory is used.
\end{proof}
